\documentstyle[%


righttag%


,11pt%


,amssymb%


,russian%               remove in Eng.


,izvru%                 Set izven% in Eng.


]{amsart}





%





\newcommand{\lft}{\langle}


\newcommand{\rgt}{\rangle}


\newcommand{\px}{\pi_X}


\newcommand{\xa}{x_{\alpha}}


\newcommand{\xak}{x_{\alpha_k}}


\newcommand{\xakk}{x_{\alpha_{k+1}}}


\newcommand{\ak}{\alpha_k}


\newcommand{\akk}{\alpha_{k+1}}


\newcommand{\dk}{\delta_k}


\newcommand{\dkk}{\delta_{k-1}}


\newcommand{\eps}{\varepsilon}


\newcommand{\Fa}{F_k(\ov x_k)+\alpha_k x_k}


\newcommand{\Fbb}{F_k(\ov x_k)+\alpha_{k+1} x_{k+1}}


\newcommand{\Fb}{F_{k-1}(\ov x_{k-1})+\alpha_k x_k}


\newcommand{\DD}{\|x_{k+1}-x_{\alpha_k}\|}


\newcommand{\D}{\|x_k-x_{\alpha_k}\|}


\newcommand{\DDD}{\|x_{k+1}-x_{\alpha_{k+1}}\|}


\newcommand{\Ax}{\|\ov x_k-x_{k+1}\|}


\newcommand{\Bx}{\|x_k-\ov x_k\|}


\newcommand{\Cx}{\|x_k-\ov x_{k-1}\|}


\newcommand{\K}{(1+\|x_*\|)}


\newcommand{\tts}[1]{}





\theoremstyle{definition}


\theorembodyfont{\rm}


\newtheorem{notenonum}{\hskip\parindent Замечание}


\renewcommand{\thenotenonum}{}





%\hoffset=-15mm%                  For Eng.


\setcounter{page}{70}


\god{2004}


\nomer{1~(500)} %                        Remove in Eng.


%% remove ^^^^^^


%\nomer{Vol.\,48, No.\,1}%               Set in Eng.


%\str{\pageref{begin}--\pageref{end}}%  Set in Eng.


\udk{517.988:519.85}





\author{\it Л.Д.\,Попов}


% NB:   ^^ remove in Eng.


\title{О схемах формирования ведущей последовательности 


в регуляризованном экстраградиентном методе


решения вариационных неравенств}


\thanks{Работа выполнена при финансовой поддержке Российского фонда


фундаментальных исследований (коды проектов \No\No\,НШ-792.2003.1,


01-01-00563).}





\date{}


%\pagestyle{myheadings}%         Set in Eng.


\pagestyle{plain} %              Remove in Eng.


\begin{document}


%\thispagestyle{plain}%          Set in Eng.


\maketitle


\label{begin}




















\section{Введение}





Многие задачи математической физики, исследования операций,


математической экономики и др. могут быть записаны в форме


вариационных неравенств, для численного решения которых к


настоящему времени предложено и исследовано большое количество


методов \cite{kinder}--\cite{kokur}, в частности, методов


градиентного типа. Как хорошо известно, для сходимости наиболее


простых из градиентных процедур необходимо, чтобы оператор,


формирующий задачу, был потенциальным (что отвечает обычным


оптимизационным постановкам) или удовлетворял усиленным свойствам


монотонности \cite{vasil}, \cite{polak}. В последующих


исследованиях ослабление условий сходимости шло главным образом по


двум направлениям, одно из которых было связано с включением в


исходную задачу регуляризирующих компонентов \cite{baku2},


\cite{baku3}, а другое опиралось на идею экстраполяции направления


спуска \cite{ant2}, \cite{korpel}, что требовало параллельного


отслеживания двух последовательностей приближений к искомому


решению: основной и экстраполирующей (называемой также прогнозной


или ведущей). В методах первого направления, получивших название


методов итеративной регуляризации, успешно решались не только


вопросы сходимости при слабых начальных предположениях


относительно свойств монотонности исходного оператора, но и


проблемы, возникающие в связи с некорректностью обсуждаемых задач


и неточностью задания их исходных данных. Платой за это выступала


необходимость постепенного и согласованного уменьшения (в пределе


до нуля) от итерации к итерации значений двух основных параметров


метода: шагового параметра и параметра регуляризации, что в


целом предопределило невысокую скорость сходимости соответствующих


алгоритмов. Второе направление привело к так называемым


экстраградиентным методам решения вариационных неравенств. Для


этих методов характерно использование хотя и малых, но конечных


значений параметра шага. Однако чувствительность экстраградиентных


методов к погрешностям вычислений существенно выше по сравнению с


методами, применяющими регуляризацию исходной задачи. Наконец,


в \cite{ant1} был опубликован метод, сочетающий в


себе идеи метода итеративной регуляризации с экстраградиентным


методoм определения направления спуска и тем самым объединивший


лучшие стороны каждого из изученных выше подходов.





Способ формирования экстраполирующей (прогнозной)


последовательности в экстраградиентных методах в том виде, как он


был впервые предложен в \cite{korpel}, основан на повторном


(двойном) вычислении направления спуска: один раз это направление


вычисляется в точке основной последовательности и один раз --- в


прогнозной точке. Однако этот способ не является единственно


возможным: при определенных условиях \cite{ppf1}, \cite{ppf2}


можно ограничиться однократным вычислением такого направления,


которое затем используется дважды: один раз при пересчете основной


точки, и второй раз --- при пересчете прогнозной точки.


Естественно, что однократное определение направления спуска


существенно снижает вычислительные затраты на проведение одной


итерации метода. 





Как будет показано в первой части данной работы,


аналогичные приемы можно сочетать с методом итеративной


регуляризации, причем вычислительная эффективность результирующего


метода только возрастает.





Вторая часть работы содержит дополнительное исследование свойств


изучаемых методов в случае, когда исходная задача не имеет


решения (является несобственной \cite{erem}). А именно, будет


показано, что в этом случае экстраградиентный метод регуляризации


доставляет некоторое квазирешение противоречивой постановки и


вариант ее трансформации в разрешимую задачу за счет минимальных


(в определенном смысле) поправок, вносимых в ее исходные данные.


Исследование этого случая следует общему подходу \cite{erem} и


опирается на результаты работ \cite{kokur}, \cite{ppf3}.





Наконец, в завершающей части работы представлены предварительные


результаты вычислительного эксперимента, подтверждающие


эффективность предлагаемых автором схем формирования прогнозной


последовательности. \vspace{1ex}








\section{Постановка задачи и описание алгоритма}





Пусть заданы $X$ --- выпуклое замкнутое множество евклидова


пространства $E^n$ и $F(x)$ --- оператор, действующий из $X$ в


$E^n$. Рассмотрим задачу решения вариационного неравенства,


которая состоит в поиске точки $x$, удовлетворяющей условиям


\begin{equation}\label{n1}


x\in X,\quad \lft F(x),\,y-x \rgt\geq0\quad


\forall y\in X.


\end{equation}


Оператор $F$ будем считать монотонным, что означает выполнение


неравенств


$$


%% \begin{equation}\label{n98}


\lft F(x)-F(y),\,x-y\rgt \geq0\qquad\forall x,y\in E^n,


%% \end{equation}


$$


и удовлетворяющим условию Липшица с константой $L$:


$$


%% \begin{equation}\label{n97}


\| F(x)-F(y)\| \leq L\,\|x-y\| \quad\forall x,y\in E^n.


%% \end{equation}


$$


Наряду с исходным неравенством (\ref{n1}) выпишем его


регуляризированный вариант: для заданного значения параметра


регуляризации $\alpha>0$ найти такую точку $\xa$, что


\begin{equation}\label{n2}


\xa\in X,\quad \lft F(\xa)+\alpha\xa,y-\xa \rgt\geq0\quad


\forall y\in X.


\end{equation}


В отличие от исходной задачи регуляризированная постановка всегда


имеет решение \cite{baioki}, и в случае разрешимости исходной


задачи выполнены соотношения \cite{baku2}, \cite{baku3}


\begin{equation}\label{n4}


\|\xa\|\leq\| x_*\|,\quad


\lim_{\alpha\to+0}\|\xa- x_*\|=0,\quad


\|\xa-x_{\gamma}\|\leq \frac{|\alpha-\gamma|}{\alpha}


\| x_*\|\ \ \forall \alpha,\gamma>0,


\end{equation}


где $x_*$ --- нормальное (т.\,е. минимальное по норме) решение


(\ref{n1}). Если же исходная постановка (\ref{n1}) неразрешима (в


других терминах


--- несобственная), то $\|\xa\|\to\infty$ при $\alpha\to+0$, однако


\begin{equation}\label{n5}


\alpha\|\xa\|\leq K_0,\quad


\lim_{\alpha\to+0} \|\alpha\,x_{\alpha}-\ov h\|=0,\quad


\|\xa-x_{\gamma}\|\leq \frac{|\alpha-\gamma|}{\alpha\gamma}


\,K_0\quad \forall \alpha,\gamma>0,


\end{equation}


где $K_0$ --- некоторая константа, $\ov h$ --- минимальный по


норме вектор среди всех таких векторов $h$, для которых


скорректированная задача (\ref{n1}) нахождения $x$ из условий


\begin{equation}\label{n6}


x\in X,\quad \lft F(x)+ h,\ y-x \rgt\geq0\quad


\forall y\in X


\end{equation}


разрешима \cite{kokur}, \cite{ppf3}. Задача (\ref{n6}) при $h=\ov


h$ может рассматриваться в качестве аппроксимирующей для


несобственной (неразрешимой в обычном смысле) исходной задачи, а


ее решение --- как некоторое обобщенное решение для (\ref{n1}).








Обратимся к основной цели исследования. Итерационный


процесс, свойства которого предлагаем рассмотреть, определяется


следующими соотношениями:


\begin{equation}\label{n3}


x_{k+1}=\px(x_k-\mu(\Fa)),\quad


\ov x_{k+1}=\px(x_{k+1}-\mu(\Fbb));


\end{equation}


здесь $\px$ --- операция проектирования на множество $X$, $\{x_k\}$


--- основная последовательность, $\{\ov x_k\}$ ---


последовательность прогнозных точек, $\mu>0$ и $\alpha_k>0$


--- числовые параметры,


\begin{equation}\label{n10}


\|F_k(x)-F(x)\|\leq\delta_k(1+\|x\|),\quad \delta_k>0,


\quad k=0,1,\dotsc


\end{equation}


Начальные приближения $x_0$, $\ov x_0$ выбираются в $X$


произвольно, например равными друг другу.





Выписанный процесс сочетает в себе идеи экстраградиентного спуска


и итеративной регуляризации. Но в отличие от своего прототипа из


\cite{ant1} метод (\ref{n3}), (\ref{n10}) на каждом шаге


вычисляет значение оператора $F$ только один раз, хотя использует


это значение дважды: при переопределении как точки $x_k$ основной


последовательности, так и точки $\ov x_k$ последовательности,


осуществляющей экстраполяцию направления спуска. Второстепенные


отличия связаны с выбором точек, в которых задача подвергается


регуляризации, и с тем, что параметр шага фиксирован на протяжении


всего хода вычислений; последнее обстоятельство, впрочем,


несущественно и нацелено лишь на упрощение последующих выкладок.








\section{Анализ сходимости в разрешимом случае}








Пусть задача (\ref{n1}) разрешима и $x_*$ --- ее минимальное по


норме решение. Анализ сходимости итерационного процесса (\ref{n14})


будем проводить


при стандартных условиях согласования значений его параметров, а


именно, будем предполагать, что


\begin{gather}\label{n14}


0<\ak,\quad 0<\dk,\quad


\lim_{k\to\infty}\ak=\lim_{k\to\infty}\dk=0,\quad


0<\mu<\frac{1}{L(1+\sqrt{2})}, \\


%


\label{n16}


\lim_{k\to\infty}\frac{\dk+\dkk}{\ak}=0,\quad


\lim_{k\to\infty}\frac{\ak-\alpha_{k+1}}{\ak^2}=0,\quad


\sum\limits_{k=1}^{\infty}\ak=\infty.


\end{gather}





В качестве примера подходящих последовательностей $\{\ak\}$ и


$\{\dk\}$ приведем следующие:


$$


\ak=(1+k)^{-\sigma},\quad \dk=(1+k)^{-\nu};\quad


\text{здесь}\ \ \


\sigma>0,\ \ \nu>0,\ \ \sigma<\min\{


\tfrac{1}{2},\,\nu\,\}.


$$








\begin{notenonum}


Условие расходимости ряда


$\sum\limits_{k=1}^{ \infty}\ak$ играет существенную роль


и потому фигурирует в общем перечне (\ref{n14}), (\ref{n16}), хотя


не является независимым, а вытекает из предшествующих ему условий.


В самом деле, сходимость последовательности $\ak$ к нулю позволяет


для любого натурального $m$ определить номер $k(m)\geq m$ такой,


что $\alpha_{k(m)}=\max\limits_{k\geq m}\,\{\ak\}$. Поскольку


$\lim\limits_{k\to\infty}\dfrac{\ak-\alpha_{k+1}}{\ak^2}=0$,


то для достаточно больших $m$ имеем


$\ak-\alpha_{k+1}<\ak^2$ \; $\forall k>m$, откуда


$$


\alpha_{k+1}>\ak(1-\ak)\geq(1-\alpha_{k(m)})\ak


\quad\forall k>m.


$$


Суммируя эти неравенства по $k$ от $m$ до произвольного $N>k(m)$,


получаем


$$


\alpha_{N+1}+\sum\limits_{k=m}^N\ak\ -\alpha_m


> (1-\alpha_{k(m)})


\sum\limits_{k=m}^N\ak,


$$


т.\,е.


$$


\sum\limits_{k=m}^{\infty}\ak  >


\sum\limits_{k=k(m)}^N\ak > 1


- \frac{\alpha_{N+1}}{\alpha_{k(m)}}


\to 1\quad \text{при} \ \ N\to\infty.


$$


Последнее соотношение, ввиду произвольности $m$, противоречит


известному свойству сходящихся рядов --- стремлению к нулю остатка


ряда.


\end{notenonum}








Сформулируем основной результат о сходимости предлагаемого


алгоритма.





\begin{thm}


Пусть задача $(\ref{n1})$ разрешима и


параметры метода $(\ref{n3}),\ (\ref{n10})$ связаны условиями


$(\ref{n14})$, $(\ref{n16})$. Тогда


\begin{equation}\label{n99}


\lim_{k\to\infty}\|x_k-x_*\|=


\lim_{k\to\infty}\|\ov x_k-x_*\|=0,


\end{equation}


где $x_*$ --- нормальное $($минимальное по норме$)$


решение задачи $(\ref{n1})$.


\end{thm}





\begin{pf}


В виду свойств (\ref{n4})


достаточно показать, что последовательность $\{x_k\}$,


генерируемая соотношениями (\ref{n3}), (\ref{n10}), обладает


свойством


$$


\lim_{k\to\infty}\|x_k-\xak\|=0,


$$


где $\xak$ --- решение регуляризированного неравенства


(\ref{n2}) при $\alpha=\alpha_k$.





Как и в \cite{ant1}, выпишем очевидное


тождество


\begin{multline*}


\D^2=\Bx^2+\Ax^2+\DD^2


+2\lft x_k-\ov x_k,\ov x_k-x_{k+1} \rgt +\\


+2\lft x_k-x_{k+1}, x_{k+1}-\xak\rgt,


\end{multline*}


которое преобразуем к виду


\begin{multline}\label{n18}


\DD^2


= \D^2-\Ax^2-\Bx^2 - \\


 -2\lft x_{k+1}-x_k+\mu(\Fa),\, \xak- x_{k+1}\rgt-\\


 -2\lft\ov x_k-x_k+\mu(\Fb),\, x_{k+1}-\ov x_k \rgt+\\


 + 2\mu\lft \Fa,\xak- x_{k+1} \rgt+


2\mu\lft \Fb,x_{k+1}-\ov x_k \rgt.


\end{multline}


Далее, пользуясь характеристическим свойством проекции на выпуклое


множество (\cite{vasil}, с.\,183)


$$


\lft\px(x)-x,y-\px(x)\rgt\geq0\quad


\forall y\in X, \ \ x\in E^n,


$$


перепишем соотношения (\ref{n3}) в виде


\begin{alignat*}{2}


\lft x_{k+1}-x_k+\mu (\Fa),\,y_1-x_{k+1}


\rgt&\geq0 &\quad &\forall y_1\in X,\\


\lft \ov x_k-x_k+\mu (\Fb),\,y_2-\ov x_k


\rgt &\geq0 &\quad &\forall y_2\in X,


\end{alignat*}


где во втором соотношении индекс $k$ заменен на $k-1$. Положим в


выписанных соотношениях $y_1=\xak$ и $y_2=x_{k+1}$. Получим


\begin{align*}


\lft x_{k+1}-x_k+\mu (\Fa),\,\xak-x_{k+1}


\rgt&\geq0,\\


\lft \ov x_k-x_k+\mu (\Fb),\,x_{k+1}-\ov x_k


\rgt&\geq0.


\end{align*}


Применим эти неравенства к правой части тождества (\ref{n18}):


\begin{multline}\label{n7}


\DD^2 \leq \D^2-\Ax^2-\Bx^2 +\\


+ 2\mu\lft \Fa,\xak- x_{k+1} \rgt +


2\mu\lft \Fb,x_{k+1}-\ov x_k \rgt=\\


= \D^2-\Ax^2-\Bx^2 + \\


+2\mu\lft \Fa,\xak-\ov x_k\rgt +


2\mu \lft F_k(\ov x_k)-F_{k-1}(\ov x_{k-1}),\,


\ov x_k - x_{k+1} \rgt.


\end{multline}


Оценим отдельно скалярные произведения, стоящие в правой части


(\ref{n7}). Начнем со второго из них как с более простого. Опираясь


на неравенство Коши--Буняковского, липшицевость оператора $F$ и


условие (\ref{n10}), получим


\begin{multline}\label{n8} 


2\mu \lft F_k(\ov x_k)-F_{k-1}(\ov x_{k-1}) ,


\ov x_k-x_{k+1} \rgt =\\


= 2\mu \lft F_k(\ov x_k)-F(\ov x_k),\,


\ov x_k-x_{k+1} \rgt+ 2\mu \lft F(\ov x_k)-F(x_k),\,


\ov x_k-x_{k+1} \rgt\ +\\


+ 2\mu \lft F(x_k)-F(\ov x_{k-1}),\,


\ov x_k-x_{k+1} \rgt


+ 2\mu \lft F(\ov x_{k-1})- F_{k-1}(\ov x_{k-1}),\,


\ov x_k-x_{k+1} \rgt\leq\\


\leq


2\mu\dk(1+\|\ov x_k\|)\Ax + 2\mu L(\sqrt{m_1}


\Bx) \frac{\Ax}{\sqrt{m_1}} + \\


+ 2\mu L(\sqrt{m_2}\Cx)\frac{\Ax}{\sqrt{m_2}}


+ 2\mu \dkk(1+\|\ov x_{k-1}\|)\Ax.


\end{multline}


Здесь $m_1,m_2>0$ --- параметры, значения которых определим


позже. А пока, используя очевидные неравенства


\begin{equation}\label{n12}


\begin{split}


\|\ov x_k\|&\leq \| \ov x_k-x_k \| + \D+\|\xak\| \leq 


\| \ov x_k-x_k \| + \D+\|x_*\|,\\


\|\ov x_{k-1}\| &\leq \Cx+\D+\|\xak\| \leq 


\Cx+\D+\|x_*\|


\end{split}


\end{equation}


и элементарное неравенство $2|ab|\leq a^2+b^2$, приводим оценку


(\ref{n8}) к виду


\begin{multline}\label{n9} 


2\mu \lft F_k(\ov x_k)-F_{k-1}(\ov x_{k-1}),\,


\ov x_k-x_{k+1} \rgt\leq\\


\leq 


\mu\dk\left(\K^2+\Bx^2+\D^2+3\Ax^2\,\right) +\\


+ \mu L\left(m_1\Bx^2+m_2\Cx^2+\left(m_1^{-1}+ m_2^{-1}\right)


\Ax^2\right) +\\


+ \mu\dkk\left(\K^2+\Cx^2+\D^2+3\Ax^2\right).


\end{multline}


Теперь вернемся к первому скалярному произведению. В силу


монотонности оператора $F$ и ограничения на точность его


вычисления, а также в силу определения (\ref{n2}) точек $\xak$


имеем


\begin{multline*}


2\mu \lft \Fa, \xak -\ov x_k \rgt=\\


= 2\mu \lft F_k(\ov x_k)-F(\ov x_k),\,\xak -\ov x_k \rgt +


 2\mu \lft F(\ov x_k)+\alpha_k x_k,\xak -\ov x_k \rgt \leq \\


\leq 2\mu \dk(1+\|\ov x_k\|)\,\|\xak -\ov x_k \| +


2\mu \lft F( \xak)+\ak\xak,\xak -\ov x_k \rgt +


2\ak\mu \lft \xak-x_k, \ov x_k-\xak\rgt \leq\\


\leq 2\mu \dk(1+\|\ov x_k\|)\|


\xak -\ov x_k \| + 2\ak\mu


\lft \xak-x_k, \ov x_k-\xak\rgt .


\end{multline*}


Первое слагаемое в правой части этой оценки удовлетворяет


соотношению


\begin{align*}


2\mu \dk(1+\|\ov x_k\|)\, \|\xak -\ov x_k \|


&\leq 2\mu \dk (1+\|\ov x_k-x_k\|+\D+\|x_*\|)(\|\xak -x_k\|+


\|x_k-\ov x_k\|)\leq\\


&\leq \mu \dk\left[2\K^2+


5\Bx^2 + 5 \D^2\ \right],


\end{align*}


что вытекает из первого неравенства (\ref{n12}), неравенства


треугольника и все того же элементарного неравенства $2|ab|\leq


a^2+b^2$. Для второго слагаемого имеем


\begin{multline*}


2\ak\mu \lft \xak-x_k, \ov x_k-\xak\rgt=


-2\ak\mu\D^2+2\ak\mu\lft \xak-x_k,\ov x_k- x_k\rgt\leq \\


\leq -2\ak\mu\D^2+2\ak\mu\D\,\Bx


\leq -\ak\mu\D^2+\ak\mu\Bx^2.


\end{multline*}


В итоге


\begin{multline}\label{n13} 


2\mu \lft \Fa, \xak -\ov x_k \rgt\leq


\mu \dk\left[2\K^2+ 5\Bx^2 +


5\D^2\ \right]- \\


- \ak\mu\D^2+\ak\mu\Bx^2.


\end{multline}


Объединим теперь соотношения (\ref{n7}) и (\ref{n9}),


(\ref{n13}). Получим


\begin{equation}\label{n15}


\DD^2\leq A_k\D^2-B_k\Ax^2-C_k\Bx^2+D_k\Cx^2+E_k,


\end{equation}


где


\begin{gather*}


A_k=1-\alpha_k\mu+\mu(6\dk+\dkk),\quad


B_k=1-\mu\,L(m_1^{-1}+m_2^{-1})-3\mu(\dk+\dkk),\\


C_k=1-\mu\,Lm_1-\alpha_k\mu-6\mu\dk,\quad


D_k=\mu(Lm_2+\dkk),\\


E_k=\mu(3\dk+\dkk)\K^2.


\end{gather*}


Теперь отойдем от стандартных


выкладок \cite{ant1}, \cite{ppf1} и свяжем величины $\DDD^2$ и


$\DD^2$ при помощи элементарного векторного неравенства


\begin{gather*}


\|x+y\|^2\geq(\|x\|-\|y\|)^2=\|x\|^2+\|y\|^2


-2\|x\|\,\|y\|


 \geq \|x\|^2+\|y\|^2-\eps\|x\|^2


-\frac{\|y\|^2}{\eps}=\\


=(1-\eps)\|x\|^2+\frac{(\eps-1)}{\eps}\|y\|^2,


\end{gather*}


которое верно для любых векторов $x$, $y$ и скаляра $\eps>0$.


Полагая здесь $x=x_{k+1}-x_{\alpha_{k+1}}$,


$y=x_{\alpha_{k+1}}-\xak$, получим


$$


\DD^2\geq (1-\eps_k)\DDD^2 +


\frac{(\eps_k-1)}{\eps_k}\|x_{\alpha_{k+1}}


-\xak\|^2.


$$


Пусть $\eps_k=\frac{ 1}{ 2}\ak\mu$, $k=0,1,\dotsc$, тогда в


силу правил согласования (\ref{n14}), (\ref{n16}) при больших $k$


имеем $1-\eps_k>0$. С учетом этого из соотношений (\ref{n4}),


(\ref{n15}) выводим заключительное рекурсивное неравенство


\begin{multline}\label{n17}


\DDD^2+\frac{ D_{k+1}}{ A_{k+1}}\Ax^2\leq


\frac{ A_k}{ 1-\eps_k}


\left( \D^2+\frac{ D_k}{ A_k}\Cx^2 \right) -\\


- \left( \frac{ B_k}{ 1-\eps_k}-


\frac{ D_{k+1}}{ A_{k+1}}\right)


\Ax^2-\frac{ C_k}{ 1-\eps_k}\Bx^2+


\frac{ E_k}{ 1-\eps_k}+ 


\frac{ (\ak-\akk)^2}{ \eps_k\ak^2}\|x_*\|^2,


\end{multline}


которое используем для обоснования сходимости метода. Начиная


с этого места, будем считать $m_1=1+\sqrt{2}$, $m_2=1$. Заметим,


что это обеспечивает легко проверяемое равенство


$m_1=m_1^{-1}+m_2^{-1}+m_2=1+\sqrt{2}$.





Рассмотрим подробнее коэффициенты в соотношении (\ref{n17}).


Во-первых,


$$


A_k=1-\ak\mu+o(\ak),\quad C_k=1-\mu L


(1+\sqrt{2})+O(\ak),


$$


в силу чего правила согласования параметров (\ref{n14}),


(\ref{n16}) обеспечивают положительность этих коэффициентов при


всех достаточно больших $k$. Во-вторых, в силу тех же причин


положительной будет разность


$$


\frac{ B_k}{ 1-\eps_k}- \frac{ D_{k+1}}{ A_{k+1}}


= \frac{ A_{k+1}B_k-(1-\eps_k)D_{k+1}}{ (1-\eps_k)A_{k+1}}.


$$


Действительно, как ранее убедились, знаменатель дроби справа


положителен, а положительность числителя при достаточно больших


$k$ вытекает из разложения


$$


A_{k+1}B_k-(1-\eps_k)D_{k+1} = 1 - \mu L(m_1^{-1}


+m_2^{-1} +m_2) +O(\ak) = 1 - \mu L(1+\sqrt{2})+O(\ak)


$$


и правил согласования параметров (\ref{n14}), (\ref{n16}).





Сказанное позволяет оставить в правой части соотношения


(\ref{n17}) только первое и два последних слагаемых, сохранив


первоначальный знак неравенства. Результат можно записать кратко в


виде


\begin{equation}\label{n19}


0\leq \omega_{k+1}\leq (1-\theta_k)\omega_k


+\sigma_k\quad\text{при всех } \ k>N,


\end{equation}


где $N$ --- достаточно большое число, и использованы обозначения


$$


\omega_k=\D^2+\frac{ D_k}{ A_k}\Cx^2 ,\quad


\theta_k= \frac{ 1-A_k-\eps_k}{ 1-\eps_k},\quad


\sigma_k= \frac{ E_k}{ 1-\eps_k}+


\frac{ (\ak-\akk)^2}{ \eps_k\ak^2}\|x_*\|^2.


$$


В силу (\ref{n14}), (\ref{n16}) при достаточно больших $k$ имеем


$\mu(6\dk+\dkk)<\frac{1}{4}\,\ak\mu$, а значит,


$$


\theta_k=\frac{ 1-A_k-\eps_k}{ 1-\eps_k} >


1-A_k-\eps_k=\frac{1}{2}\ak\mu-\mu(6\dk+\dkk)>


\frac{1}{4}\,\ak\mu.


$$


Это вместе с предположением о расходимости ряда


$\sum\limits_{k=1}^{\infty} \ak$ влечет расходимость ряда


$\sum\limits_{k=1}^{\infty} \theta_k$. Кроме того, из тех же


условий следует


$$


0\leq\frac{\sigma_k}{\theta_k} \leq 4\left(


\frac{\mu(3\dk+\dkk)(1+\|x_*\|)^2}{(1-\eps_k)\ak\mu}+


\frac{(\ak-\akk)^2}{\eps_k\ak^2\ak\mu}\|x_*\|^2\right)


\to0 \quad \text{при } \ k\to\infty.


$$


Остается сослаться на известный результат (\cite{vasil}, с.\,90),


согласно которому последовательность $\{\omega_k\}$,


удовлетворяющая (\ref{n19}) при условии расходимости


ряда $\sum\limits_{k=1}^{\infty} \theta_k$ и сходимости к нулю


отношения $r_k=\frac{\sigma_k}{\theta_k}$, сходится к


нулю. Вместе с этой последовательностью сходятся


к нулю и последовательности $\|x_k-\xak\|$, $\|x_k-\ov


x_{k-1}\|$, что с учетом соотношений (\ref{n4}) влечет искомые


равенства (\ref{n99}).


\end{pf}








\section{Анализ несобственного случая}





Пусть теперь задача (\ref{n1}) неразрешима. Рассмотрим


параметрическое семейство возмущенных задач отыскания точек $x$,


удовлетворяющих условиям


\begin{equation}\label{n80}


x\in X,\quad \lft F(x)+h,\, y-x \rgt \geq0\quad


\forall y\in X;


\end{equation}


здесь $h$ --- параметр, описывающий постоянное возмущение при


вычислении значения оператора~$F$. Как известно, множество $H$


всех таких значений параметра $h$, при которых возмущенная задача


(\ref{n80}) имеет хотя бы одно решение, не пусто и почти выпукло


\cite{kokur}, \cite{ppf2}. Последнее означает, что его замыкание


$\ov H$ выпукло и удовлетворяет условию


$$


\bold{ri}\,\ov


H\subseteq H\subseteq \ov H,


$$


где $\bold{ri}$ обозначает


операцию взятия относительной внутренности выпуклого множества.





Среди возмущенных задач вида (\ref{n80}), обладающих решением,


естественно попытаться выбрать задачу, наиболее близкую к


исходной, и считать ее обычное решение некоторым обобщенным


решением (или квазирешением) исходной неразрешимой (несобственной)


постановки. Близость исходной и возмущенной задачи можно измерять


нормой вектора $h$, например, евклидовой нормой $\|h\|$.





Обозначим $\|\ov h\|=\min\limits_{h\in\ov H} \|h\|$. Такой элемент


очевидно существует и является единственным в силу выпуклости и


замкнутости $\ov H$. Хотя, строго говоря, $\ov h$ может не


лежать в $H$, свойство почти выпуклости множества $H$ гарантирует


существование в нем элементов, сколь угодно близких к $\ov h$.


Это служит хорошей мотивировкой для изучения задачи поиска $\ov


h$. Следующая теорема устанавливает связь между этой задачей и


экстраградиентным методом регуляризации.





\begin{thm}


Пусть задача $(\ref{n1})$ неразрешима и $\xa$ ---


решение регуляризированной задачи $(\ref{n2})$, отвечающее


некоторому $\alpha>0$. Тогда


\begin{itemize}


\item[(a)] $\alpha\xa\in H$,


\item[(b)] $\xa$ является решением возмущенной


задачи $(\ref{n80})$ при $h=\alpha\xa$,


\item[(c)] $\lim\limits_{\alpha\to+0}


\|\alpha\xa-\ov h\|=0$,


\end{itemize}


где $\ov h$ --- минимальный по норме элемент множества $\ov H$.





Если $($в дополнение к предыдущему$)$


параметры метода $(\ref{n3})$, $(\ref{n10})$


связаны условиями $(\ref{n14})$, $(\ref{n16})$ и значения


$\ak$ не возрастают, то также


\begin{itemize}


\item[(d)] $\lim\limits_{k\to\infty}


\|\ak x_k-\ov h\|=0$,


\end{itemize}


т.е. последовательность $\{\ak x_k\}$ сходится к


оптимальному параметру коррекции $\ov h$.


\end{thm}





\begin{pf}


Первые три пункта теоремы доказаны


в \cite{kokur}, \cite{ppf3}. Обоснуем последний ее


пункт. С этой целью проследим последовательность рассуждений,


выполненную при выводе неравенства (\ref{n17}) в доказательстве


теоремы 1. Предположение о неразрешимости исходной задачи меняет


только два звена этой последовательности. Первое звено связано с


оценкой $\|\xak\|\leq\|x_*\|$, которая благодаря (\ref{n5}) теперь


приобретает вид $\|\xak\|\leq K_0\ak^{-1}$. Соответственно


следует заменить $\|x_*\|$ на $K_0\ak^{-1}$ в соотношениях


(\ref{n12})--(\ref{n13}) и в их следствиях. Второе


звено касается оценки


$$


\|\xakk-\xak\|\leq \frac{ (\ak-\akk)}{ \ak}\|x_*\|,


$$


которая применялась в (\ref{n17}) и теперь имеет вид


$$


\|\xakk-\xak\|\leq \frac{ (\ak-\akk)}{ \ak\akk}K_0.


$$


С учетом этих изменений перепишем соотношение (\ref{n17}):


\begin{multline}\label{n17a}


\DDD^2+\frac{ D_{k+1}}{ A_{k+1}}\Ax^2\leq


\frac{ A_k}{ 1-\eps_k}


\left( \D^2+\frac{ D_k}{ A_k}\Cx^2 \right) -\\


- \left( \frac{ B_k}{ 1-\eps_k}-


\frac{ D_{k+1}}{ A_{k+1}}\right)


\Ax^2-\frac{ C_k}{ 1-\eps_k}\Bx^2+


\frac{ \mu(3\dk+\dkk)(1+K_0\ak^{-1})^2}{ 1-\eps_k}+ \\


+\frac{ (\ak-\akk)^2}{ \eps_k\ak^2\,\akk^2} K_0^2.


\end{multline}


Умножим левую и правую части выписанного неравенства на $\akk^2$.


Оставляя справа только неотрицательные слагаемые и опираясь на то,


что значения $\ak$ не возрастают, получим аналог оценки


(\ref{n19}):


\begin{multline}\label{n17b}


0\leq\omega'_{k+1}\leq


\frac{ A_k}{ 1-\eps_k}\ \omega'_k+


\frac{ \mu(3\dk+\dkk)(K_0+\ak)^2}{ 1-\eps_k}


+ \frac{ (\ak-\akk)^2}{ \eps_k\ak^2}K_0^2=\\


=(1-\theta_k)\,\omega'_k+\sigma'_k\quad\forall k>N,


\end{multline}


где введены обозначения


$$


\omega'_k= \|\ak x_k-\ak\xak\|^2


+\ak^2\frac{ D_k}{ A_k}\Cx^2,\ \


\sigma'_k=


\frac{ \mu(3\dk+\dkk)(K_0+\ak)^2}{ 1-\eps_k}\


+ \frac{ (\ak-\akk)^2}{ \eps_k\ak^2}K_0^2,


$$


и где $N$ --- номер, начиная с которого второе и третье слагаемые


в правой части неравенства (\ref{n17a}) можно считать


неположительными.





Нетрудно убедиться, что условия согласования параметров итерационного


процесса (\ref{n3}) при всех достаточно больших $k$ обеспечивают неравенство


$$


0\leq r'_k=\frac{\sigma'_k}{\theta_k}<


\frac{\sigma'_k}{ 1-A_k-\eps_k}=


\frac{ \sigma'_k}{ O(\ak)}


\to0\quad\text{при } \ k\to\infty.


$$


Кроме того те же условия, как было показано ранее, обеспечивают


расходимость ряда $\sum\limits_{k=1}^{\infty}\theta_k$, что


позволяет вновь применить, но уже к последовательности


$\{\omega'_k\}$, утверждение (\cite{vasil}, с.\,90), в


силу которого эта последовательность сходится к нулю. По


определению вместе с этой последовательностью сходится к нулю и


последовательность $\|\ak x_k-\ak\xak\|$, что с учетом соотношений


(\ref{n5}) завершает доказательство последнего пункта теоремы~2.


\end{pf}








\section{Результаты вычислительного эксперимента}





Вычислительный эксперимент проводился на линейных задачах о


дополнительности с кососимметричной матрицей коэффициентов.


Напомним, что для заданной квадратной матрицы $A$ линейная задача


о дополнительности состоит в поиске пары векторов $x$, $y$,


удовлетворяющих условиям


$$


y=Ax+b,\quad x\geq0,\ \ y\geq0,\quad \lft x,y\rgt=0.


$$


Линейные задачи о дополнительности с кососимметричными матрицами,


т.\,е. матрицами, удовлетворяющими условию $A=-A^T$, являются


наиболее простым примером вариационных неравенств \cite{kinder},


\cite{png}, чей оператор $F(x)=Ax+b$ удовлетворяет условию


монотонности в наиболее слабой форме


$$


\lft F(x)-F(y),\ x-y \rgt=\lft A(x-y),


x-y\rgt=0\quad\forall x,y.


$$





Все тестовые примеры имели заранее известное решение. Погрешности


вычисления моделировались датчиком случайных чисел с равномерным


распределением. Величина шагового параметра выбиралась таким


образом, чтобы обеспечить наибольшую скорость сходимости каждого


из алгоритмов при фиксированных стратегиях выбора значений других


параметров. Эти стратегии определялись соотношениями


$$


\ak=\alpha_0(1+k)^{-1/2},\quad \dk=\delta_0(1+k)^{-1},


\quad k=1,2,\dotsc


$$





В эксперименте участвовали три метода: метод итеративной


регуляризации из \cite{baku2}, экстраградиентный метод


регуляризации из \cite{ant1} и предлагаемый в данной работе


модифицированный экстраградиентный метод с альтернативной схемой


формирования прогнозной последовательности. Во всех случаях


начальное приближение к искомому решению выбиралось одинаковым.


Типичные результаты расчетов представлены ниже в таблице, где в


крайней левой колонке помещен номер итерации, а в остальных


--- показатели достигнутой к этому моменту точности (т.\,е.


отклонение текущего приближения от решения тестовой задачи) в


логарифмической шкале.


\medskip








\begin{center}


{\bf Точность, достигаемая разными методами


за одинаковое число шагов}


\smallskip





\begin{tabular}{|c|c|c|c|}


\hline


\No & Метод итеративной & Экстраградиентный метод & Модифицированный \\


итерации & регуляризации & регуляризации & экстраградиентный метод \\


\hline %% \hline


0 & 1.36956 & 1.36956 & 1.36956 \\


3000 & $-0.99844$ & $-2.13194$ & $-2.05402$ \\


6000 & $-1.62098$ & $-2.71738$ & $-2.46219$ \\


9000 & $-1.62454$ & $-3.30334$ & $-2.86882$ \\


12000 & $-1.64115$ & $-3.89139$ & $-3.27424$ \\


15000 & $-1.62102$ & $-4.45678$ & $-3.67904$ \\


18000 & $-1.63893$ & $-4.99480$ & $-4.08608$ \\


21000 & $-1.64986$ & $-5.72244$ & $-4.50322$ \\


36000 & $-1.67753$ & $-5.91552$ & $-5.83104$ \\


\hline


\end{tabular}


\end{center}


\medskip





Из приведенных данных видно, что оба экстраградиентных метода


регуляризации превосходят по скорости сходимости метод итеративной


регуляризации и вполне сопоставимы в этом плане друг с другом.


Вместе с тем, затрачивая на достижение одной и той же точности


решения примерно одинаковое число итераций, эти два метода


разнятся по общему времени вычислений, т.\,к. итерация второго из


них выполняется на компьютере примерно в два раза быстрее, чем у


первого. Таким образом, предварительные результаты тестирования


показали существенное преимущество новой схемы формирования


прогнозной последовательности в экстраградиентном методе


регуляризации.








\begin{thebibliography}{99}








\bibitem{kinder} Киндерлерер Д., Стампаккья Г. {\em Введение в


вариационные неравенства и их приложения}. -- М.: Мир, 1983.


-- 256~с.





\bibitem{glow} Гловински Р., Лионс Ж.Л., Тремольер Р. 


{\em Численное исследование вариационных неравенств}. -- М.: Мир,


1979. -- 574~с.





\bibitem{baioki} Байокки К., Капело А.


{\em Вариационные и квазивариационные неравенства}. -- М.: Наука,


1988. -- 488~с.





\bibitem{png} Harker P.T., Pang J.-S.


{\em Finite-dimensional variational inequalities and


nonlinear complementarity problems$:$ a survey of


theory, algorithms and applications} //


Math. Program. -- 1990. -- V.\,48. --


P.\,161--220.





\bibitem{t-a} Тихонов А.Н., Арсенин В.Я. {\em Методы решения


некорректных задач}. -- М.: Наука, 1974. -- 223~с.





\bibitem{vasil} Васильев Ф.П. {\em Методы оптимизации}. -- М.:


Факториал, 2002. -- 823~с.





\bibitem{t-l-ya} Тихонов А.Н., Леонов А.С., Ягола А.Г.


{\em  Нелинейные некорректные задачи}. -- М.: Наука, 1995.


-- 307~с.





\bibitem{polak} Бакушинский А.Б., Поляк Б.Т. {\em О решении


вариационных неравенств} // ДАН СССР. -- 1974. --


Т.\,219. -- \No\,5. -- С.\,1038--1041.





\bibitem{baku2} Бакушинский А.Б., Гончарский А.В. {\em Итеративные


методы решения некорректных задач}. -- М.: Наука,


1989. -- 127~с.





\bibitem{baku3} Бакушинский А.Б., Гончарский А.В. {\em Некорректные


задачи. Численные методы и приложения}. --


М.: Изд-во МГУ, 1989. -- 197~с.





\bibitem{ant2} Антипин А.С. {\em Градиентный и экстраградиентный


подходы в билинейном равновесном программировании}. --


М.: Изд-во ВЦ РАН, 2002. -- 86~с.





\bibitem{ant1}  Антипин А.С., Васильев Ф.П. {\em Регуляризованный


экстраградиентный метод для решения вариационных неравенств}


// Вычисл. методы и программиров. -- 2002. -- Т.\,3.


-- \No\,2. -- С.\,144--150.





\bibitem{kokur} Кокурин М.Ю. {\em Операторная регуляризация и


исследование нелинейных монотонных задач}. --


Йошкар-Ола: Изд-во Марийск. ун-та, 1998. -- 292~с.





\bibitem{korpel} Корпелевич Г.М. {\em Экстраградиентный метод для отыскания


седловых точек и других задач} // Экономика и матем. методы.


-- 1976. -- Т.\,12. -- Вып.\,4. -- С.\,747--756.





\bibitem{ppf1} Попов Л.Д. {\em Модификация метода Эрроу--Гурвица поиска


седловых точек} // Матем. заметки. -- 1980. -- Т.\,28.


-- Вып.\,5. -- С.\,777--784.





\bibitem{ppf2} Попов Л.Д. {\em Об одной модификации метода


Эрроу--Гурвица поиска седловых точек с адаптивной процедурой


определения итерационного шага} / В кн. Классификация и


оптимизация в задачах управления. -- Свердловск: Уральск. НЦ АН СССР.


-- 1981. -- С.\,52--56.





\bibitem{erem} Еремин И.И., Мазуров Вл.Д., Астафьев Н.Н.


{\em Несобственные задачи линейного и выпуклого программирования}.


-- М.: Наука, 1983. -- 336~с.





\bibitem{ppf3} Попов Л.Д. {\em О применении метода проекции для


нахождения аппроксимационных корней монотонных отображений} //


Изв. вузов. Математика. -- 1995. -- \No\,12.


-- С.\,74--80.





\end{thebibliography}





\vskip 2cm





\post{\it Институт математики и}{\it Поступила}


\post{\it механики Уральского отделения}{07.08.2003}


\post{\it Российской Академии наук}{}





\label{end}


\end{document}


